\emph{Conventions.} After the Wick rotation (Fact \claimref{synthesis.imported-one-loop-integrals-and-pauli-villars-regulators}(ii)) all momenta are Euclidean, $Q(q) > 0$, and the Dirac algebra is the Clifford algebra of the Euclidean vectors, $\gamma_{\alpha}\gamma_{\beta} + \gamma_{\beta}\gamma_{\alpha} = 2\delta_{\alpha\beta}$. The Feynman rules of \claimref{field-theory.coupled-system-tree-level} give, with the units collected outside the product as in \claimref{synthesis.one-loop-qed-in-the-spacetime-algebra} (the field theory's unit $I_{4}$ acting as the right unit, numbers read in $\mathrm{span}\{1, I_{4}\}$): fermion propagator $-k/Q(k)$ with one unit (numerator the vector $k$), vertex $-e\gamma_{\alpha}$ with one unit, photon $(\delta_{\alpha\beta} - (1-\xi)r_{\alpha}r_{\beta}/Q(r))/Q(r)$ (the covariant form, which needs the Krein-form premise \claimref{synthesis.krein-form-premise-indefinite-inner-product-for-gauge-fields}), and $-1$ per closed fermion loop. Per loop $e^{2}\int d^{d}k/(2\pi)^{d} = a(4\pi)^{\epsilon}\int d^{d}k/\pi^{d/2}$, $a = \alpha/(4\pi)$. The overall factors are $(-1)(-eI_{4})^{2}(-I_{4})^{2} = -e^{2}$ at one loop and $(-1)(-eI_{4})^{4}(-I_{4})^{4} = -e^{4}$ at two loops. The trace of a product of vectors is $4\langle\cdot\rangle_{0}$ (\claimref{synthesis.one-loop-qed-in-the-spacetime-algebra}(1)); no unit enters a vector-like loop, so the units are coefficients and the unit result's obstruction to continuing $I_{4}$ does not arise. Dimensional regularisation (Fact \claimref{synthesis.imported-dimensional-regularisation-and-the-bmhv-gamma5}(i)) supplies the continuation in $d$ with poles at $d = 4$ and its Lorentz and vector-like gauge invariance. The following standard rules are its clause (iv) (\cite{ChetyrkinTkachov1981}; Fact \claimref{synthesis.imported-dimensional-regularisation-and-the-bmhv-gamma5}(iv)): $\int d^{d}\ell/\pi^{d/2}\,Q(\ell)^{r}(Q(\ell) + \Delta)^{-N} = \Gamma(r + d/2)\Gamma(N - r - d/2)\Delta^{d/2 + r - N}/(\Gamma(d/2)\Gamma(N))$; symmetric integration $\ell^{\alpha_{1}}\cdots\ell^{\alpha_{2r}} \to Q(\ell)^{r}\{\delta^{r}\}/(d(d+2)\cdots(d+2r-2))$; translation invariance; the vanishing of scale-free integrals; and $\int d^{d}k\,\partial_{k}\cdot(\cdots) = 0$. The gauge parameter is kept general; since the result does not depend on it, it holds in the Feynman gauge, whose Fock space realises the one-particle Gupta-Bleuler form, restated in the store as the trace form of Clifford conjugation (\claimref{synthesis.clifford-conjugation-is-the-twistor-adjoint-and-its-trace-form-is-the-one-particle-gupta-bleuler-form}; a restatement, not a derivation of (K1)).

\emph{Step 1 (the algebra in $d$ dimensions).} For a vector $b$ and a product $X$ of vectors, $b\gamma_{\alpha} = -\gamma_{\alpha}b + 2b_{\alpha}$ gives $\gamma_{\alpha}Xb\gamma_{\alpha} = -(\gamma_{\alpha}X\gamma_{\alpha})b + 2bX$, and $\gamma_{\alpha}\gamma_{\alpha} = d$; by induction on the length of $X$ every contraction $\gamma_{\alpha}X\gamma_{\alpha}$ is a polynomial in $d$ times products of the vectors of $X$. Both rules hold in $\Cl(D, 0)$ and, with the reciprocal frame, in $\Cl(1, D-1)$ for every integer $D$; the check compares them with the explicit sums in $\Cl(D, 0)$, $D = 4, \ldots, 8$, and $\Cl(1, D-1)$, $D = 4, 5, 6$, on words of length up to four (control: $d = 4$ in place of $D = 6$ fails). The trace of a product of vectors follows from $\mathrm{tr}(a_{1}\cdots a_{n}) = \sum_{j \ge 2}(-1)^{j}(a_{1}\cdot a_{j})\,\mathrm{tr}(\cdots\hat a_{1}\cdots\hat a_{j}\cdots)$, $\mathrm{tr}\,1 = 4$, which uses only the anticommutator. The resulting traces are polynomials of degree at most two in $d$; the check compares every word below (Feynman and gauge parts, both projections) with $4\langle\cdot\rangle_{0}$ of the explicit products summed over a basis in $\Cl(D, 0)$, $D = 4, 5, 6$: agreement at three values of $d$ for a polynomial of degree two in $d$ is an identity in $d$, for the sampled momenta.

\emph{Step 2 (the diagrams).} The fermion loop carries $k$ and $k - q$; the internal photon carries $r = k - l$. With $\Pi_{\mu\nu}^{(2)} = -e^{4}\int\int \frac{d^{d}k\,d^{d}l}{(2\pi)^{2d}}\,\sum_{j}T_{j}$:
\[
T_{\mathrm{se}1} = \frac{\mathrm{tr}[\gamma_{\mu}k\gamma_{\alpha}l\gamma_{\beta}k\gamma_{\nu}(k-q)]P_{\alpha\beta}}{Q(k)^{2}Q(k-q)Q(l)Q(r)},\quad
T_{\mathrm{se}2} = \frac{\mathrm{tr}[\gamma_{\mu}k\gamma_{\nu}(k-q)\gamma_{\alpha}(l-q)\gamma_{\beta}(k-q)]P_{\alpha\beta}}{Q(k)Q(k-q)^{2}Q(l-q)Q(r)},
\]
\[
T_{\mathrm{vx}} = \frac{\mathrm{tr}[\gamma_{\mu}k\gamma_{\alpha}l\gamma_{\nu}(l-q)\gamma_{\beta}(k-q)]P_{\alpha\beta}}{Q(k)Q(k-q)Q(l)Q(l-q)Q(r)},\qquad P_{\alpha\beta} = \delta_{\alpha\beta} - (1-\xi)\frac{r_{\alpha}r_{\beta}}{Q(r)},
\]
the matrix order following the fermion line against its arrow, with momentum conserved at every vertex. The scalar function is $\Pi = \Pi_{\mu\mu}/((d-1)Q(q))$, and transversality is $q_{\mu}q_{\nu}\Pi_{\mu\nu} = 0$. At one loop $\Pi^{(1)}_{\mu\nu} = -e^{2}\int\mathrm{tr}[\gamma_{\mu}k\gamma_{\nu}(k-q)]/(Q(k)Q(k-q))$. By homogeneity $\Pi^{(n)} \propto Q(q)^{-n\epsilon}$, so one may set $Q(q) = 1$.

\emph{Step 3 (reduction).} With $D_{1} = Q(k)$, $D_{2} = Q(k-q)$, $D_{3} = Q(l)$, $D_{4} = Q(l-q)$, $D_{5} = Q(k-l)$ and $Q(q) = 1$: $k\cdot q = (D_{1} - D_{2} + 1)/2$, $l\cdot q = (D_{3} - D_{4} + 1)/2$, $k\cdot l = (D_{1} + D_{3} - D_{5})/2$, so every numerator is a polynomial in the $D_{i}$ and every term an integral $I(a_{1}, \ldots, a_{5}) = \int\int\prod D_{i}^{-a_{i}}$ (measure $d^{d}k\,d^{d}l/\pi^{d}$) with integer $a_{i}$. (a) $a_{5} \le 0$: expand $D_{5}^{n} = (D_{1} + D_{3} - 2k\cdot l)^{n}$; the $k$ integral of $(k\cdot l)^{j}$ over $D_{1}^{-A}D_{2}^{-B}$ is, by Feynman parameters and the integrals of the Conventions paragraph,
\[
\sum_{r}\binom{j}{2r}\frac{(2r-1)!!}{2^{r}}\,Q(l)^{r}(l\cdot q)^{j-2r}\,\frac{\Gamma(A + B - r - \tfrac d2)\Gamma(j - r + \tfrac d2 - A)\Gamma(r + \tfrac d2 - B)}{\Gamma(A)\Gamma(B)\Gamma(j + d - A - B)},
\]
zero if $A \le 0$ or $B \le 0$ (scale-free), and the $l$ integral is a bubble. (b) $a_{1} \le 0$ (or $a_{2} \le 0$; or, by $k \leftrightarrow l$, $a_{3}$ or $a_{4} \le 0$): shift $k = k' + q$ (or not), expand the numerator in $Q(k')$ and $k'\cdot q$, use the same formula with external momentum $P = l - q$ (or $l$), which leaves $Q(P)^{d/2 + r - A - B}$, and the $l$ integral is a bubble with one non-integer index. (c) All $a_{i} > 0$: integration by parts with $\partial_{k}\cdot(k - l)$ gives the triangle rule
\[
(d - a_{1} - a_{2} - 2a_{5})\,I = a_{1}\mathbf{1}^{+}(\mathbf{5}^{-} - \mathbf{3}^{-})I + a_{2}\mathbf{2}^{+}(\mathbf{5}^{-} - \mathbf{4}^{-})I,
\]
since $2k\cdot(k - l) = D_{1} + D_{5} - D_{3}$, $2(k - q)\cdot(k - l) = D_{2} + D_{5} - D_{4}$; each step lowers $a_{3} + a_{4} + a_{5}$ by one, so the recursion ends in (a) or (b). The bubble (Fact \claimref{synthesis.imported-dimensional-regularisation-and-the-bmhv-gamma5}(iv)) is $G(\alpha, \beta) = \Gamma(\alpha + \beta - \tfrac d2)\Gamma(\tfrac d2 - \alpha)\Gamma(\tfrac d2 - \beta)/(\Gamma(\alpha)\Gamma(\beta)\Gamma(d - \alpha - \beta))$ \cite{ChetyrkinTkachov1981}. In (a) the Gamma functions have arguments $n \pm d/2$ and $n + d$ in the same numbers as in $M_{1} = G(1,1)^{2}$; in (b) as in $M_{2} = G(1,1)G(1, 3 - d/2)$ (arguments $n - d$ and $n + 3d/2$ once each); integer shifts of an argument $z$ give rational factors by $\Gamma(z + 1) = z\Gamma(z)$, so every $I$ is $R_{1}(d)M_{1} + R_{2}(d)M_{2}$ with $R_{1}, R_{2}$ rational. Benchmarks (checked): $I(1,1,1,1,0) = M_{1}$; $I(1,1,1,1,-1) = -\tfrac12M_{1}$ (from $\int k^{\mu}/(D_{1}D_{2}) = \tfrac12q^{\mu}G(1,1)$); $I(0,1,1,1,1) = M_{2}$; and $I(1,1,1,1,1) = 2(3 - d)M_{1}/(d - 4) + R_{2}M_{2} = 6\zeta_{3} + O(\epsilon)$.

\emph{Step 4 (expansion).} $G(1,1) = g(\epsilon)/\epsilon$ with $g = \Gamma(1+\epsilon)\Gamma(1-\epsilon)^{2}/((1-2\epsilon)\Gamma(1-2\epsilon))$, and
\[
\frac{M_{2}}{M_{1}} = \frac{G(1, 1+\epsilon)}{G(1,1)} = \frac{(1 - 2\epsilon)\,\Gamma(1+2\epsilon)\Gamma(1-2\epsilon)^{2}}{2(1 - 3\epsilon)\,\Gamma(1+\epsilon)^{2}\Gamma(1-\epsilon)\Gamma(1-3\epsilon)} = \frac{1 - 2\epsilon}{2(1 - 3\epsilon)}\exp\Bigl(-6\zeta_{3}\epsilon^{3} - 9\zeta_{4}\epsilon^{4} + \cdots\Bigr)
\]
by $\ln\Gamma(1 + z) = -\gamma z + \sum_{n \ge 2}(-1)^{n}\zeta_{n}z^{n}/n$ \cite{AbramowitzStegun1964}: the coefficients of $\gamma$ and $\zeta_{2}$ are $2 - 4 - 2 + 1 + 3 = 0$ and $\tfrac12(4 + 8 - 2 - 1 - 9) = 0$. The rational functions are expanded exactly at $d = 4 - 2\epsilon$. The $\overline{\mathrm{MS}}$ normalisation replaces $(4\pi)^{\epsilon}$ per loop by $e^{\gamma\epsilon}$, and $e^{\gamma\epsilon}g = 1 + 2\epsilon + O(\epsilon^{2})$.

\emph{Step 5 (one loop).} $\mathrm{tr}[\gamma_{\mu}k\gamma_{\mu}(k - q)] = 4(2 - d)\,k\cdot(k - q)$ and $k\cdot(k - q) = (D_{1} + D_{2} - 1)/2$, whose first two terms are scale free; so $\Pi^{(1)} = -\tfrac{2(d-2)}{d-1}a(4\pi)^{\epsilon}G(1,1)Q(q)^{-\epsilon} = a(4\pi)^{\epsilon}Q(q)^{-\epsilon}\,\tfrac{g(\epsilon)}{\epsilon}\bigl(-\tfrac43 + \tfrac49\epsilon + O(\epsilon^{2})\bigr)$; the $\overline{\mathrm{MS}}$ constant is $\tfrac49 + 2(-\tfrac43) = -\tfrac{20}{9}$. With $e^{2}_{\mathrm{eff}}(q) = e_{0}^{2}/(1 - \Pi(Q(q)))$ (the Euclidean propagator $\sum_{n}Q(q)^{-1}\Pi^{n}$), $1/a_{\mathrm{eff}} = 1/a_{0} - \Pi^{(1)}/a_{0} - \Pi^{(2)}/a_{0} + \cdots$, and $d(1/a_{\mathrm{eff}})/d\ln Q(q) = -\beta_{0} - \beta_{1}a$: the one-loop term gives $\beta_{0} = \tfrac43$, screening.

\emph{Step 6 (two loops).} The check computes the three traces (Step 1), reduces them (Step 3) and expands (Step 4). Results, exact in $d$: the coefficient of $(1 - \xi)$, summed over the three diagrams, is zero as a rational function of $d$ in both masters (it is not zero for the vertex diagram alone); $q_{\mu}q_{\nu}\Pi^{(2)}_{\mu\nu}$ is zero in both masters, in the Feynman part and the gauge part. With $\Pi^{(2)} = a^{2}(4\pi)^{2\epsilon}Q(q)^{-2\epsilon}(R_{1}M_{1} + R_{2}M_{2})$, the series $R_{1} + R_{2}M_{2}/M_{1}$ begins at $\epsilon^{1}$ with $-2\epsilon + (16\zeta_{3} - \tfrac{31}{3})\epsilon^{2}$: there is no pole of order three or two, $\Pi^{(2)} = a^{2}(-2/\epsilon + \cdots)$, and the $\overline{\mathrm{MS}}$ constant is $16\zeta_{3} - \tfrac{31}{3} + 4(-2) = 16\zeta_{3} - \tfrac{55}{3}$. In the Feynman gauge the two self-energy insertions together give $+\tfrac43a^{2}/\epsilon^{2}$ and the vertex correction $-\tfrac43a^{2}/\epsilon^{2}$: the one-loop subdivergences, removed by the counterterms $Z_{2} - 1$ on the self-energy subgraphs and $Z_{1} - 1$ on the vertex subgraphs, cancel in the sum because $Z_{1} = Z_{2}$ at one loop (\claimref{synthesis.one-loop-qed-in-the-spacetime-algebra}(4)), and the bare sum needs only the overall counterterm. In $1/a_{\mathrm{eff}}$ the term $-\Pi^{(2)}/a_{0} = -a(-2/\epsilon)Q(q)^{-2\epsilon} + \cdots$ has $\ln Q(q)$ coefficient $-4a$, so $\beta_{1} = 4$ and $Z_{3} = 1 - \tfrac43 a/\epsilon - 2a^{2}/\epsilon$; with $n_{f}$ fermions each diagram has one closed loop and both coefficients scale with $n_{f}$. In $\alpha$: $\mu\,d\alpha/d\mu = 2(4\pi)[\beta_{0}(\alpha/4\pi)^{2} + \beta_{1}(\alpha/4\pi)^{3}] = \tfrac{2}{3\pi}\alpha^{2} + \tfrac{1}{2\pi^{2}}\alpha^{3}$. Comparison (computed values against the literature, read 2026-11-02): Steinhauser's $Z_{3}$ at order $\alpha\alpha_{s}$ with $C_{F} \to 1$, $\alpha_{s} \to \alpha$ is $1 - (\alpha/\pi)/(3\epsilon) - (\alpha/\pi)^{2}/(8\epsilon) = 1 - \tfrac43a/\epsilon - 2a^{2}/\epsilon$ \cite{Steinhauser2002}; the constant $f_{2} = \tfrac{55}{12} - 4\zeta_{3}$ of \cite{FanchiottiKniehlSirlin1993}, with the one-loop $f_{1} = \tfrac53$ and the relative factor $\tfrac34\alpha/\pi = 3a$, gives $-\tfrac43a\cdot 3a\,(\tfrac{55}{12} - 4\zeta_{3}) = a^{2}(16\zeta_{3} - \tfrac{55}{3})$, as computed. The coefficient $\beta_{1}$ is that of \cite{JostLuttinger1950} (Fact \claimref{synthesis.imported-lsz-reduction-ward-identity-and-the-second-qed-beta-coefficient}(iii)).

\emph{Step 7 (Ward identity at two loops).} In $\Cl(1,3)$, $S(k) = (k - m)^{-1} = (k + m)/(Q(k) - m^{2})$ with the units of the Feynman rules outside the product, as the right unit $I_{4}$ (\claimref{synthesis.one-loop-qed-in-the-spacetime-algebra}(5)), and $S(k + h) - S(k) = S(k+h)[(k - m) - (k + h - m)]S(k) = -S(k + h)hS(k)$. For an open line $C(p) = V_{0}S(p_{1})V_{1}\cdots S(p_{n})V_{n}$ (the momenta $p_{j}$ differ from $p$ by loop momenta), let $T_{j}$ be the line with $p_{1}, \ldots, p_{j-1}$ shifted by $h$ and the rest not; the insertion of the vertex $h$ at line $j$ is $T$ with $S(p_{j})$ replaced by $S(p_{j} + h)hS(p_{j}) = -(S(p_{j} + h) - S(p_{j}))$, which is $-(T_{j+1} - T_{j})$, so the sum over $j$ is $-(C(p + h) - C(p))$. The two-loop self-energies have open lines of three propagators (rainbow, crossed) and one propagator (the line of the polarisation insertion), and the closed loop of the insertion telescopes in the same way to the difference of one integrand at loop momenta $l + h$ and $l$, which integrates to zero by translation invariance. Hence $h_{\mu}\Lambda^{(2)\mu}(p + h, p) = -(\Sigma^{(2)}(p + h) - \Sigma^{(2)}(p))$; to first order in $h$, $\frac{d}{dt}S(k + th)|_{0} = h/(Q(k) - m^{2}) - 2(k\cdot h)(k + m)/(Q(k) - m^{2})^{2} = -S(k)hS(k)$, so $\Lambda^{(2)\mu}(p, p) = -\partial\Sigma^{(2)}/\partial p_{\mu}$, and on shell $Z_{1} = Z_{2}$ at two loops \cite{Ward1950}. The check verifies the three open-line identities, the closed-loop identity and the derivative identity exactly at rational momenta with $m = \tfrac32$, with two controls (the sum without the middle insertion, and $\gamma_{0}$ in place of $h$, both fail). Proved for the integrands; the passage to the regulated integrals uses the shift invariance of dimensional regularisation, and the vertex and self-energy diagrams are the standard ones (argued).

\emph{Status of the clauses.} (1) to (4) are exact computations given the imported rules of dimensional regularisation and the Gamma-function expansion (computed); (5) is proved for the integrands and argued for the integrals. Transport-absent; the spacetime algebra is convenient (the continuation in $d$ is one Clifford rule checked in every integer dimension from four to eight). The regulator is imported, not native.
